% !TeX root = ../papers/nc-p3.tex
\section{座標1の単元剰余を表す楕円曲線と三次までの全水準}
\label{sec:rank-fifteen-vertical-unit-levels}
\begin{lemma}\label{lem:rank-fifteen-vertical-elliptic-residues}
標数3の $\bar x=1$, $\bar y\ne0$ の剰余表示は
\[
 \bar y=(v-v^3)^2,\qquad
 k(v)\bigl(\sqrt v,\sqrt{v-1},\sqrt{v+1}\bigr)
\]
から得られる。ただし $k$ は代数閉体とする。この関数体の滑らかな完備曲線は
種数1で、$\bar y$ 直線への写像は次数48のGalois被覆である。
\end{lemma}
\begin{proof}
$K=1/v$, $U=K^2$, $H=K-v$, $V=H^2$ とすれば $V(U+3)/(U-1)^2=1$,
$V^3/(U-1)^4=(v-v^3)^2$ となる。$\bar y\ne0$ なら $v\ne0,\pm1$ である。
$i^2=-1$, $w^2=v$ とし、$A+B=\sqrt{1+v}$, $A-B=\sqrt{1-v}$ で $A,B$ を取る。
$A^2+B^2=1$, $2AB=v$ より $E=A^2-B^2$ は $E^2=1-v^2$ を満たす。
$n=EA/v$ とすると $H/n=-EB/v$,
$n^2+(H/n)^2=K^2-1$ となる。他の四根は順に
\[
 P_+=(v+1)/w,\quad P_-=i(A+B),\quad Q_+=(v-1)/w,\quad Q_-=A-B
\]
と取れる。いずれも $v\ne0,\pm1$ で使える。

$v,v-1,v+1$ の平方類は三つの相異なる単純零点での付値により独立である。
従って $v$ 直線への被覆は次数8で、$0,1,-1,\infty$ にのみ分岐する。
各点では単元は完備局所体で平方なので慣性指数は2である。
Riemann--Hurwitz公式から $2g-2=-16+4\cdot4=0$ となる。
一方 $v\mapsto\epsilon v+b$ ($\epsilon=\pm1$, $b\in\mathbb F_3$) は
$(v-v^3)^2$ を保つ六変換である。三つの被開平式を符号を除き置換するため
いずれも次数8の被覆へ持ち上がり、全次数48だけの自己同型を与える。
具体的には $A=(1-t^2)/(1+t^2)$, $B=2t/(1+t^2)$ と置けるので、
$W^2=t-t^3$ の表示では $w=W/(1+t^2)$ となる。
\end{proof}

\begin{theorem}\label{thm:rank-fifteen-vertical-unit-levels}
標数零の通常な階数15窓で、3の上のある付値で $\bar x=1$ とし、
$\bar y\ne0$ の $\mathbb F_3$ 上の次数が高々3とする。この窓は標準Koszul左右AS正則で、
各始点で四中心四次式上自由階数256である。
\end{theorem}
\begin{proof}
次数1は定理\ref{thm:rank-fifteen-unit-residues}による。
\path{computations/verify_rank_fifteen_vertical_unit_levels.py} は素体上の
モニックな二次・三次式を列挙し、既約な三つの二次式と八つの三次式の
それぞれについて有限体表示を検証する。必要な拡大次数は $6,8,12,18$ である。
指定した $\bar y$ の最小多項式、$v=1/K$, $H=K-v$ の二式と全ての平方根を照合する。
各表示で未切断経路の長さ0から5の次元は $(1,4,10,20,35,56)$、
四次中心の文字別次元は $(3,2,0,0)$ である。文字0の三つの自由座標の後二つと、
文字1の二つの自由座標を選び、四始点全てで長さ13までの中心商の上界を得る。

所定の $x,y$ へはまず $(U+3)^3=(x^3/y)(U-1)^2$ の指定した根を持ち上げる。
$U-1$ は単元であり、剰余で微分は非零である。
$V=x(U-1)^2/(U+3)$ も単元で、残る根はモニックな式から順に持ち上がる。
単元主元、一般の経路と中心の下界、中心商の上界の一致は
定理\ref{thm:rank-fifteen-unit-residues}と同じ全次数の証明を与える。
有限体の共役によりこれら十一式の全ての根に適用できる。
\end{proof}
\begin{corollary}
$x=1$, $y\ne0$ で $[\mathbb Q(y):\mathbb Q]\le3$ とする。
全ての添字保持同型類は標準Koszul左右AS正則で、
系\ref{cor:rank-fifteen-vertical-one-rational}の $t$ の式による四類
（$y=1/27$ のみ二類）であり、最小周期はすべて2である。
より一般に $y$ の原始整数最小多項式の斉次化が、3で還元して次数3以下の
閉点を零点に持てば同じ結論が成り立つ。
\end{corollary}
\begin{proof}
3での正付値と負付値は前節で尽くされた。全共役が単元ならば、
所定の剰余閉点の上の共役根を選び、本定理を適用する。次数が高々3なら、
還元した原始斉次式の零点のいずれも次数が高々3である。
\end{proof}
